\section{Related Work}
\label{sec:related}

Our paper sits at the intersection of three literature threads.
Cross-corpus emotion analysis has documented label-space
incompatibility without prescribing adaptation methods.
Adversarial domain adaptation has been validated on sentiment and
image tasks but not on cross-dataset emotion.
Class-imbalance regularization (focal loss) has been applied to
emotion classification but not combined with adversarial frameworks
on the cross-dataset benchmark we target.

\paragraph{Cross-corpus emotion analysis.}
\citet{bostan2018analysis} systematically catalogued the label-space
incompatibility between emotion datasets, demonstrating that
projecting heterogeneous corpora to a shared Ekman scheme requires
either dropping incompatible classes or accepting contested mappings
(e.g.\ ISEAR's \emph{shame} and \emph{guilt} have no consensus
projection). Their analysis stops at diagnosis: they do not propose
methods to adapt models across the resulting heterogeneous splits.
\citet{demszky2020goemotions} released GoEmotions and reported a
source-only cross-dataset transfer baseline to ISEAR and EmoInt, but
did not test domain adaptation methods against this baseline.
A concurrent line of work proposes unified emotion benchmarks
spanning multiple corpora: \citet{koufakou2024towards} consolidate
eleven publicly-available English single-labeled datasets into a
single benchmark and compare unified-corpus training against
per-corpus baselines, finding that unified training improves some
target corpora but degrades others. Their evaluation, however, is
restricted to a mixed-training protocol against single-corpus
baselines; they do not adopt a leave-one-dataset-out (LODO)
protocol, they do not compare domain-adaptation methods, and they do
not investigate class-imbalance regularization. \citet{plaza2024emotion}
review the state of emotion analysis in NLP and explicitly identify
the lack of systematic cross-corpus method comparison as an open
challenge.
Our work picks up where these four leave off: we adopt the Ekman-6
projection (excluding shame/guilt per Bostan and Klinger's
caveat), evaluate under both mixed-training and LODO protocols, and
test six methods that span the standard cross-domain answer
(DANN/CDAN) and the class-imbalance answer (focal loss). The LODO
protocol exposes a regime that mixed-training cannot — held-out target
generalization — which is where our headline finding emerges.

\paragraph{Adversarial domain adaptation in NLP.}
DANN \citep{ganin2016dann} introduced the gradient reversal layer
and sigmoid $\lambda$ annealing, with primary validation on
cross-domain sentiment (Amazon product reviews) and image tasks.
CDAN \citep{long2018cdan} extends DANN by conditioning the
discriminator on task-head class predictions via a multilinear map,
demonstrated on image domain adaptation (Office-31, ImageCLEF).
The multi-domain text classification extension
\citep{chen2018multinomial} generalizes DANN to more than two source
domains, again on sentiment. Cross-dataset \emph{emotion} (as opposed
to sentiment or topic) has not been systematically evaluated under
these methods in the published literature to the best of our
knowledge. Our negative results on adversarial methods for this task
(Section~\ref{sec:results}) connect to the survey-level caveat in
\citet{wilson2020dasurvey} that domain adaptation does not
universally transfer across task families.

\paragraph{Class imbalance and focal loss.}
Focal loss \citep{lin2017focal} reweights the cross-entropy loss to
down-weight well-classified examples and concentrate gradient on hard
or rare examples; it was originally proposed for dense object
detection but has been adapted to imbalanced NLP classification
\citep{suresh2021labelaware}. In emotion classification specifically,
class imbalance is pronounced: GoEmotions is dominated by joy;
ISEAR has no surprise; WASSA-21 essay distributions skew toward
sadness (see Table~\ref{tab:datastats}). Recent work has applied
focal loss to imbalanced emotion benchmarks in single-corpus or
cross-lingual settings: \citet{vazquez2024pcicunam} combine focal
loss with weighted cross-entropy and a hierarchical classifier on a
five-language emotion-detection shared task, demonstrating gains over
plain cross-entropy on individual languages. Their setup, however,
holds the corpus constant and varies the language; our setup holds
the language constant (English) and varies the corpus, which is a
structurally different generalization problem.
The interaction between focal loss and the cross-dataset shift
addressed by adversarial methods (DANN+Focal, CDAN+Focal) has not,
to our knowledge, been published for text emotion classification at
the cross-corpus level. Our DANN+Focal and CDAN+Focal cells
(Section~\ref{sec:results}) provide the first such measurement on the
three-emotion-dataset benchmark.

\paragraph{Position of this work.}
We do not propose a new method. We provide the first systematic
empirical comparison of two complementary mitigations — adversarial
feature alignment and focal-loss class-imbalance regularization —
on a cross-dataset emotion benchmark with pre-registered protocol and
seed budgets. Our headline finding (focal alone outperforms all
adversarial variants on the one target where any method helps) is
non-obvious from the prior literature and informs the practical
choice of mitigation for related cross-corpus text classification
problems.
